\documentclass[11pt]{article}
\usepackage{geometry}
\geometry{
letterpaper,
top=1in,
left=1in,
right=1in,
bottom=1in
}

% required packages
\usepackage[hidelinks]{hyperref}
\usepackage{titlesec}
\usepackage{color}
\usepackage{verbatim}
\usepackage{enumitem}
\usepackage{tabularx}

% line spacing
\usepackage{setspace}
\setlength{\parindent}{20pt}
\setlength{\parskip}{1em}

% margin adjustments
\raggedbottom

% ====================
% DOCUMENT STARTS HERE
% ====================
\begin{document}

\onehalfspacing

Before we get to the core question, I would like to define "native language". A quick google search of the query "what is someone's native language" gives me many search results, the first two being Wikipedia and Reddit. I love the internet. Wikipedia's definition of a native language is "A first language (L1), native language, native tongue, or mother tongue is the first language a person has been exposed to from birth or within the critical period". This raises some very interesting questions for me, because the primary language I speak is English, yet that was not the first language I was exposed to, nor was it the only one I learned growing up as a child. The first language I was exposed to was Sichuanese, a dialect of Mandarin, when I was living in an orphanage for the first thirteen months of my life. Afterwards, I was adopted by my parents and they taught me English for the first five years of my life. My family and I proceeded to move to China for job reasons, and I started learning Standard Mandarin (spoken and written) while attending international school. I moved back to the States when I was around 10 years old, thus solidifying my language of choice as English. In short, if you consider my "native language" the one I was exposed to first, then I know roughly 50 words in Mandarin and 0 in Sichuanese.

Already this assignment has confounded me in ways not anticipated. However, for simplicity's sake, and one that is probably the most accurate to the definition of "native language", my native language is English. Now that that's out of the way I'm going to get back to the core question. The first thing I have done is run this exact document (up until this point right here) using the methods shown in class. (wait, no, THIS POINT). The output of that file was 159 words! So I know about 159 words confirmed. Ok, but now how do I "know" these words? Just because I can string together a coherent sentence doesn't necessarily mean that I "know" these words. Let's use some philosophy to define how I "know" a word. The branch of philosophy that deals with the theory of knowledge is epistemology which describes the nature of knowledge and the different types of knowledge. According to epistemology, knowledge is an awareness, understanding, familiarity or skill. To "have knowledge" one must make successful contact with reality. Through this contact, we can understand, interpret, and interact with the world. Of course, there are many counterarguments to what knowledge really is and how we can ever know anything, or if there is no knowledge at all. Thus, for this assignment, and before I go down too many rabbit holes of what it means to "know" something, knowing a word means that I can do one of the following: define it by recalling its dictionary definition, ascertain its general meaning using context clues from surrounding words in a sentence, or interact with it by writing it into a sentence that others can understand and interpret themselves.

Now that I have come to an abitrary philosophical way of defining what it means to know a word, let's define what an English word is morphologically.

% ====================
% DOCUMENT ENDS HERE
% ====================
\end{document}

